\documentclass{article}
\usepackage{icmlko}
\usepackage{graphicx}
\usepackage{booktabs}
\usepackage{amsmath}

\title{문턱은 빌려온 수였다 \\
  {\large 「예: 5\%」가 관문의 상수로 서 있었고, 자료에 물으니 \textbf{3.3배} 헐거웠다}}

\begin{document}
\maketitle

\begin{abstract}
이 저장소의 순열 판정은 「귀무 세계가 진짜와 나란히 놓일 수 있는가」를 먼저 묻는다 ---
\emph{비교가능성 관문}. 그 관문의 문턱은 900노트 동안 \texttt{COMPARABLE\_REL} $=0.05$ 였다.
이번 노트는 **그 수가 어디서 왔는지**를 저장소에서 찾고, 없으면 자료에서 다시 세웠다.
\claimx{🔴 \textbf{출처는 찾았다. 그런데 근거가 아니었다} ---
사전등록이 인용한 것은 이슈의 \textbf{「문턱(예: 5\%)」}이라는 \emph{예시 수} 한 줄이다.}
{게다가 같은 저장소가 \textbf{같은 축의 5\%}(기준 MAE 의 5\% $=0.519$ 일)를
「출처가 관례뿐이라 안 씀」으로 \textbf{이미 거부했다} --- 크기로 거부한 수를 관문 문턱으로 쓰고 있었다.}
\claimx{🔴🔴 \textbf{자료가 정한 문턱은 $\mathbf{0.015122}$} --- 옛 문턱의 $0.302$ 배다.}
{정의: \textbf{기준 팔 차이 하나만으로 그 뽑기의 개선이 진짜를 넘어설 수 있으면 비교 불가다.}
그러면 $T=G/(L_{\mathrm{used}}\cdot c_r)$ 이고, 관례 상수를 한 개도 안 쓴다.}
\claimx{🔴 \textbf{옛 문턱은 판정을 뒤집는 크기의 $\mathbf{3.31}$ 배를 허용하고 있었다.}}
{$0.05\times L_{\mathrm{used}}\times c_r = \mathbf{1.318}$ 일 --- \textbf{채택 크기 $1.0$ 일보다도 크다.}
관문이 지키려던 판정의 빈틈은 $G=0.399$ 일이다.}
\claimx{🔴 \textbf{그 문턱에서 직전 노트의 「한 칸 밖에서도 산다」 판이 $\mathbf{200/200}$ 비교 불가가 된다.}}
{판 ① 원판 전량은 $k=0/200$ 으로 \textbf{살고}, 판 ② 만 전량이 걸린다.
⚠ $k\cdot p$ 는 안 바뀐다 --- 그 노트는 관문으로 부분집합을 안 만들었다.}
\end{abstract}

\section{문턱은 잰 적이 없는 유일한 상수였다}

이 저장소는 자를 여러 번 고쳤다 --- 구간은 BCa 로(규약 47), 천장은 LOO 로(규율 4),
관문 신고는 $k\cdot u\cdot r$ 로(규약 48). 🔴 \textbf{그런데 문턱 그 자체는 한 번도 안 쟀다.}
사전등록은 「5\% 는 지금 박는다. 결과를 보고 고치지 않는다」라고 적었고, 그 문장은 지켜졌다 ---
\textbf{결과를 보고 안 고쳤다는 것과 근거가 있다는 것은 다른 문장이다}.

\begin{table}[h]\centering\small
\begin{tabular}{ll}
\toprule
어디 & 무엇이 적혀 있나 \\
\midrule
\texttt{prereg\_922:95} & 「근거: \textbf{\#175 가 예로 든 수}」 \\
이슈 \#175 51줄 & 「기후값 MAE 차가 \textbf{문턱(예: 5\%)} 안인지」 \\
\texttt{목표.md:138} & 「기준 MAE 의 5\% $=0.519$ 일 (\textbf{출처가 관례뿐이라 안 씀})」 \\
\bottomrule
\end{tabular}
\end{table}

두 번째 근거로 적힌 「옛 귀무는 19.319\% 로 이 문턱을 네 배 넘는다」는
\textbf{문턱이 무엇을 걸러낸 적이 있다는 사후 확인}이지 크기의 근거가 아니다 ---
19.319\% 를 걸러내는 문턱은 $0.05$ 말고도 무한히 많다.

\section{「비교 불가」를 판정에서 끌어낸다}

관문의 물음은 「이 뽑기를 판정에 쓸 자격이 있나」다. 판정은 순위 검정이므로
\textbf{자격을 잃는 지점은 순위를 뒤집을 수 있는 지점}이다.

\[
T=\frac{G}{L_{\mathrm{used}}\cdot c_r},\qquad
G=\min_i(R-y_i),\qquad
L_{\mathrm{used}}=\max\!\big(1,\ \hat{L}^{\,\mathrm{BCa}}_{\mathrm{hi}}\big)
\]

$c_r$ 은 진짜의 기준 팔 MAE, $R$ 은 진짜 개선, $y_i$ 는 $i$번째 귀무의 개선이다.
$L$ 은 \emph{전달률} --- 기준 팔이 하루 나빠지면 개선이 몇 일 움직이나.
개선 $=$ 기준 $-$ 시험 이 200뽑기 전부에서 부동소수 마지막 비트까지 성립하므로,
\textbf{시험 팔이 전혀 안 따라오면 전달률은 정확히 $1$} 이다. 그것이 하한선이고,
실측 BCa 상한이 그보다 크면 그 값을 쓴다(더 엄해진다).

\begin{table}[h]\centering\small
\begin{tabular}{lrrrr}
\toprule
팔 & $G$(일) & $c_r$(일) & $L_{\mathrm{used}}$ & $T$ \\
\midrule
① 원판 전량 \textbf{(정본)} & $0.39868$ & $12.1284$ & $2.1738$ & $\mathbf{0.015122}$ \\
② \texttt{b\_prv}$=-1$ 제외 & $0.20088$ & $10.3394$ & $1.0$ & $0.019429$ \\
진단 ③ & $4.12077$ & $53.4436$ & $1.0$ & $0.077105$ \\
\bottomrule
\end{tabular}
\end{table}

🔴 \textbf{사전등록이 미리 적은 예상은 틀렸다.} 「$T\approx0.033$」이라고 적었는데
그것은 $L=1$ 판의 값($0.032872$)이었다. 전달률의 BCa 가 $[-0.944,\,+2.174]$ 로
\textbf{판정 불능}이고, 사전등록이 \emph{미리} 「상한을 쓴다」를 박아 뒀으므로
$T$ 가 그 절반 아래로 내려갔다. \textbf{예상은 틀렸고 규칙은 지켰다.}

\begin{figure}[h]\centering
\includegraphics[width=\linewidth]{figs/fig1.pdf}
\caption{상대차 분포와 두 문턱(가로축 로그). 판 ① 은 새 문턱에서도 $k=0/200$ 로 \textbf{살고},
판 ② 와 진단 ③ 은 \textbf{전량}이 걸린다.}
\end{figure}

\section{옛 문턱이 허용하던 것}

\begin{figure}[h]\centering
\includegraphics[width=\linewidth]{figs/fig2.pdf}
\caption{(a) 문턱이 허용하던 왜곡을 \textbf{개선 눈금}으로 옮기면 $1.318$ 일 ---
판정을 뒤집는 크기 $G=0.399$ 일의 $\mathbf{3.31}$ 배이고 채택 크기 $1.0$ 일보다 크다.
(b) 팔마다 $T$ 가 $5.10$ 배 갈린다. (c) 배선 검정력 세 눈금.}
\end{figure}

\begin{table}[h]\centering\small
\begin{tabular}{lll}
\toprule
팔 & 옛 문턱 $0.05$ & 🔴 $T=0.015122$ \\
\midrule
935 ① 원판 전량 & $k{=}0/200$, $r{=}0.099$ & $k{=}\mathbf{0}/200$, $r{=}0.328$ \\
935 ② 그 칸 제외 & $k{=}0/200$, $r{=}0.982$ & 🔴 $k{=}\mathbf{200}/200$, $r{=}3.247$ \\
935 진단 ③ & $k{=}200/200$ & $k{=}200/200$, $r{=}12.93$ \\
933 & $k{=}0/200$, $r{=}\mathbf{0}$ & $k{=}0/200$, $r{=}\mathbf{0}$ \\
925 G2 & $k{=}0/\mathbf{1}$, $u{=}0.95$ & $k{=}0/\mathbf{1}$, $u{=}0.95$ \\
\bottomrule
\end{tabular}
\end{table}

관문 신고는 규약 48 재개정판 문안 그대로다 ---
\textbf{「이 $B$뽑기에서 $k$번 걸렸다 $\cdot$ 뽑기당 초과확률 95\% 상한 $u$(Clopper--Pearson)
$\cdot$ 관측 최댓값 $\div$ 문턱 $=r$」}. 「통과」와 「검정력 0」은 금지어이고,
러너가 \textbf{자기 출력에서 그 낱말을 자가검사}한다(첫 판에서 내 주석 문자열이 걸렸다).

\section{덤 --- 규약을 소급하다 부딪힌 것}

「통과」는 저장소에 \textbf{4,838줄 / 1,055파일}, 「검정력 0」은 \textbf{136줄 / 48파일} 있다.
그런데 취합 검사(⑤$'$)는 \textbf{게이트 명부를 \texttt{git grep -lF -e '"통과":'} 로 뽑는다} ---
🔴 \textbf{키 이름까지 금지하면 취합 검사가 게이트를 못 찾는다.}
금지는 \emph{관문의 결론 문장}에만 걸어야 한다(부칙).
그리고 산출물은 물론 \textbf{사전등록 문서도 안 고쳤다} --- 소급 수정은
「측정 전에 박았다」는 성질 자체를 없앤다.

\section{못 넘는 것}

🔴 \textbf{이 자는 판정 통계량에서 나온다} --- 그래서 관문이 판정과 독립이 아니다.
관문의 물음이 「이 뽑기가 판정을 뒤집을 수 있나」인 이상 판정 눈금 없이는 정의할 수 없다
(사전등록 §E-1 에 \emph{미리} 적었다).
$G$ 는 \textbf{관측된} 최소 빈틈이라 뽑기를 늘리면 줄어든다 --- $T$ 는 200뽑기에서 잰 값이다.
세 팔의 $T$ 가 $5.10$ 배 갈리므로 \textbf{상수 하나로 쓰는 것 자체가 임시}다.
이 노트는 새 팔을 안 쟀다 --- 크기 판정은 \textbf{해당 없음}이다.

\end{document}
